\section{Genetic Algorithms}

The Python part of this exercise is available in \myhref{https://github.com/LosDelDGIIM/LosDelDGIIM.github.io/blob/main/subjects/Erasmus-DUE/GKI/Exercises/Exercise4.ipynb}{this Jupyter Notebook}.\\

Imagine a delivery driver who needs to visit a set of cities and return to the starting city. The goal is to find the shortest possible route that visits each city exactly once. This is a classic optimization problem known as the Traveling Salesperson Problem (TSP).

\subsection*{Genetic Algorithm Concepts}

\begin{ejercicio}
    Explain the fundamental concepts of genetic algorithms, including population, fitness function, selection, crossover, and mutation.\\

    Genetic algorithms are a type of optimization algorithm inspired by the process of natural selection. They work by maintaining a population of candidate solutions, which evolve over time through selection, crossover, and mutation.
    \begin{itemize}
        \item \ul{Population}: A set of candidate solutions to the optimization problem. Each solution is often represented as a chromosome (e.g., a string of bits, a permutation, etc.).
        \item \ul{Fitness Function}: A function that evaluates the quality of each candidate solution. It assigns a fitness score based on how well the solution solves the problem.
        \item \ul{Selection}: The process of choosing candidate solutions from the population to create offspring for the next generation. Common selection methods include roulette wheel selection, tournament selection, and rank-based selection.
        \item \ul{Crossover}: A genetic operator that combines two parent solutions to produce one or more offspring. The crossover operator is designed to exchange information between the parents, potentially creating better solutions.
        \item \ul{Mutation}: A genetic operator that introduces random changes to a candidate solution. Mutation helps maintain genetic diversity in the population and allows the algorithm to explore new areas of the solution space.
    \end{itemize}
\end{ejercicio}
\begin{ejercicio}
    How do genetic algorithms differ from traditional optimization algorithms?

    While traditional optimization algorithms (like gradient descent, linear programming, etc.) often work with a single solution and improve it iteratively, genetic algorithms maintain a population of solutions and use stochastic processes to explore the solution space.
\end{ejercicio}
\begin{ejercicio}
    What are the advantages and disadvantages of using genetic algorithms for optimization problems?
    \begin{itemize}
        \item \ul{Advantages}:
        \begin{itemize}
            \item Can handle complex, non-linear, and multi-modal optimization problems.
            \item Does not require gradient information, making it suitable for problems where the fitness function is not differentiable.
            \item Can escape local optima due to its stochastic nature.
        \end{itemize}
        \item \ul{Disadvantages}:
        \begin{itemize}
            \item Can be computationally expensive, especially for large populations and many generations.
            \item May require careful tuning of parameters (population size, mutation rate, etc.) for good performance.
            \item No guarantee of finding the global optimum.
        \end{itemize}
    \end{itemize}
\end{ejercicio}

\subsection*{TSP Representation}
\begin{ejercicio}
    Explain how to define a fitness function for the TSP that evaluates the quality of a route.

    A suitable fitness function for the TSP can be defined as the inverse of the total distance of the route. For example, if the total distance of a route is \(D\), the fitness can be calculated as:
    \[\text{Fitness} = \frac{1}{D}
    \]
    This way, shorter routes will have higher fitness values, and the genetic algorithm will be incentivized to evolve towards shorter routes.
\end{ejercicio}

\subsection*{Genetic Operators for TSP}
\begin{ejercicio}
    Describe suitable crossover and mutation operators for the TSP. Explain why standard crossover and mutation operators might not be effective.
    
    For the TSP, standard crossover and mutation operators (like one-point crossover or bit-flip mutation) are not effective because they can produce invalid routes (e.g., duplicate cities or missing cities). Instead, specialized operators that preserve the permutation structure of the solution are needed. Examples of crossover operators suitable for TSP include:
    \begin{itemize}
        \item \ul{Variant of Ordered Crossover (OX)}: A subrange of the first parent is copied to the child (in the same positions), and the remaining positions are filled with values from the second parent. The variant considers that if a value from the second parent is already in the child, it is skipped until a non-duplicate value is found.
        
        % // TODO: Comprobar que no es exactamente, sino una variante
    \end{itemize}
    Examples of mutation operators suitable for TSP include:
    \begin{itemize}
        \item \ul{Swap Mutation}: Two cities in the route are selected and their positions are swapped.
        \item \ul{Inversion Mutation}: A segment of the route is selected and reversed.
    \end{itemize}
\end{ejercicio}
\begin{ejercicio}
    Discuss the importance of maintaining valid routes (permutations) during crossover and mutation.

    In the TSP, each candidate solution is a permutation of cities. It is crucial to maintain valid routes during crossover and mutation to ensure that each city is visited exactly once. Standard crossover and mutation operators might produce invalid routes (e.g., duplicate cities or missing cities), which would not be meaningful solutions to the TSP. Therefore, specialized operators that preserve the permutation structure are necessary to ensure the integrity of the solutions. Invalid routes would lead to incorrect fitness evaluations and hinder the convergence of the genetic algorithm towards optimal solutions.
\end{ejercicio}